\section{COMPOSE: a molecular rewrite process}
\label{sec:compose}

\method instantiates the preceding construction as a stochastic process on
two-dimensional molecular graphs. This section fixes the chemistry boundary,
rewrite basis, parameterization, and training contract. Exact model and corpus
identities are reported only after the gated editing rebuild is frozen.

\subsection{Molecular state and closure contract}

\paragraph{Broad-organic state.}
A non-null state contains at most 40 active atoms represented by an ordered
15-class element--valence vocabulary, explicit bond and aromaticity state,
implicit hydrogen counts, and formal charge. Production uses persistent slots
so an action can address the same atom across a trajectory. A slot can hold an
element, a null value, or an inert deletion marker; only the authoritative
element predicate contributes a molecular atom. Canonical molecular identity
is used to merge semantic states, but exact slot-addressed tensors are retained
for executable replay.

A frozen applicability census over 800 benchmark lead molecules admits
800/800 under the broad-organic, charge-retaining support and 264/800 under a
neutral C/N/O/F support; every benchmark lead contains at most 33 heavy atoms.
These are support-membership counts, not generation, likelihood, or
optimization results. They establish only that the declared 40-atom
broad-organic representation does not exclude this particular benchmark panel.

The process models a 2D molecular graph. Stereochemistry, isotope labels,
radicals, and charge-changing edits are outside the current editing contract.
Charged molecules may appear, but a retained charged center cannot be created,
deleted, moved, retyped, or have an incident bond changed by a training edit.
This charge-preserving rule prevents a graph edit from silently standing in for
an unspecified protonation or standardization operation.

\paragraph{Closure.}
The legal enumerator applies family-specific structural masks before a mark can
be sampled. The executor then checks the complete proposed rewrite, including
active-slot state, bond representation, per-atom valence, connectedness, and
molecular sanitization. Accepted non-null successors therefore remain in the
declared state space. This is a graph-validity and representation claim, not a
claim of synthetic accessibility, favorable three-dimensional geometry, or
biological activity.

\begin{figure}[t]
  \centering
  \begin{tikzpicture}[
  x=1cm,
  y=1cm,
  font=\sffamily\scriptsize,
  panel/.style={
    draw=black!13,
    fill=composepanel,
    rounded corners=4pt,
    line width=0.5pt
  },
  state/.style={
    draw=composelink!62,
    fill=white,
    rounded corners=3pt,
    minimum height=0.62cm,
    minimum width=0.75cm,
    inner sep=3pt,
    line width=0.6pt
  },
  actionmark/.style={
    draw=composelink!48,
    fill=composeaccent,
    rounded corners=2pt,
    minimum height=0.43cm,
    minimum width=0.57cm,
    inner sep=2pt,
    line width=0.5pt
  },
  formula/.style={
    draw=composelink!48,
    fill=white,
    rounded corners=3pt,
    align=center,
    inner xsep=5pt,
    inner ysep=4pt,
    line width=0.55pt
  },
  flow/.style={
    -{Stealth[length=4.5pt,width=3.5pt]},
    draw=composelink,
    line width=0.7pt
  },
  bidir/.style={
    {Stealth[length=4.2pt,width=3.2pt]}-{Stealth[length=4.2pt,width=3.2pt]},
    draw=composelink,
    line width=0.65pt
  },
  muted/.style={text=black!62},
  paneltitle/.style={font=\sffamily\bfseries\small, anchor=west}
]
  % Panel 1: semantic state and trans-dimensional moves.
  \node[panel, minimum width=5.00cm, minimum height=4.05cm, anchor=south west]
    (p1) at (0,0) {};
  \node[paneltitle] at (0.24,3.72) {1.\; Executable state changes};

  \node[state] (null) at (0.58,2.28) {$\varnothing$};
  \node[state] (xone) at (1.82,2.28) {$\mathcal X_{1}$};
  \node[font=\sffamily\bfseries, text=black!45] at (2.78,2.28) {$\cdots$};
  \node[state, minimum width=0.92cm] (xn) at (3.76,2.28) {$\mathcal X_{n}$};

  \draw[bidir] (null) -- node[above, muted] {\scriptsize source} (xone);
  \draw[bidir] (xone) -- node[above, muted] {\scriptsize birth/death} (2.55,2.28);
  \draw[bidir] (3.02,2.28) -- (xn);
  \draw[flow] (xn.north east)
    to[out=52,in=128,looseness=4.2]
    node[pos=0.5, below=2pt, fill=composepanel, inner sep=1pt, muted]
      {\scriptsize fixed-$n$ rewrite}
    (xn.north west);

  \node[align=center, text width=4.35cm, muted] at (2.50,0.83) {%
    Each legal mark is a complete graph rewrite.\\
    Non-null successors satisfy the declared molecular contract.};

  % Panel 2: marks execute and aliases merge by molecular identity.
  \node[panel, minimum width=5.70cm, minimum height=4.05cm, anchor=south west]
    (p2) at (5.20,0) {};
  \node[paneltitle] at (5.44,3.72) {2.\; Execute marks, quotient aliases};

  \node[state] (x) at (5.67,2.24) {$x$};
  \node[draw=black!16, fill=white, rounded corners=3pt, minimum width=1.55cm,
        minimum height=2.02cm] (fiber) at (7.14,2.13) {};
  \node[font=\sffamily\bfseries, muted] at (7.14,2.91) {$\mathcal A(x)$};
  \node[actionmark] (a1) at (7.14,2.66) {$a_1$};
  \node[actionmark] (a2) at (7.14,2.15) {$a_2$};
  \node[actionmark] (a3) at (7.14,1.64) {$a_3$};
  \node[state, minimum width=0.86cm] (y1) at (9.65,2.48) {$[y_1]$};
  \node[state, minimum width=0.86cm] (y2) at (9.65,1.55) {$[y_2]$};

  \draw[flow] (x) -- (fiber.west);
  \node[muted] at (8.52,2.96) {\scriptsize execute $T$};
  \draw[flow] (a1.east) -- (y1.west);
  \draw[flow] (a2.east) -- (y1.west);
  \draw[flow] (a3.east) -- (y2.west);

  \node[align=center, text width=5.02cm, muted] at (8.05,0.63) {%
    Distinct slot- or symmetry-level marks can execute to the same
    canonical molecular successor.};

  % Panel 3: pushforward kernel and successor-level control.
  \node[panel, minimum width=5.90cm, minimum height=4.05cm, anchor=south west]
    (p3) at (11.10,0) {};
  \node[paneltitle] at (11.34,3.72) {3.\; Molecular kernel and control};

  \node[formula, text width=4.96cm] (marklaw) at (14.05,2.96) {%
    $\lambda_\theta(a\mid x,t)
      =\Lambda_\theta(x,t)\,p_\theta(a\mid x,t)$};
  \node[formula, text width=4.96cm] (pushforward) at (14.05,2.02) {%
    $\displaystyle
      P_\theta(y\mid x,t)
      =\sum_{a:\,T(x,a)\cong y}p_\theta(a\mid x,t)$};
  \node[formula, text width=4.96cm] (doob) at (14.05,0.93) {%
    $\displaystyle
      P^g_b(y\mid x)
      =P(y\mid x)\frac{h_{b-1}(y)}{h_b(x)}$};

  \draw[flow] (marklaw) -- (pushforward);
  \draw[flow] (pushforward) -- (doob);
\end{tikzpicture}

  \caption{\textbf{Rewrite Generator Matching with executable molecular
  events.} Atom birth and death cross cardinality strata; legal marks execute
  to molecular successors; aliases are aggregated by canonical molecular
  identity; and finite-horizon control acts on the resulting successor kernel.}
  \label{fig:framework}
\end{figure}

\subsection{Minimal editing basis}
\label{sec:operator-basis}

Table~\ref{tab:operators} gives the eight-family \emph{Active8} basis entering
Gate~0, T1, and P50. It is a development-only support freeze, not final
production support. Atom birth and death alter cardinality; the other families
change composition, electronics, connectivity, or topology at fixed
cardinality. The final checkpoint must declare its exact support or the table
and every support-sensitive experiment must be revised before results are
inspected.

\begin{table}[t]
  \centering
  \caption{\textbf{Editing rewrite basis entering the gated rebuild.}
  ``Same'' means the active-atom cardinality is unchanged. Internal model-slot
  names for cycle operations are intentionally omitted.}
  \label{tab:operators}
  \small
  \begin{tabularx}{\linewidth}{@{}lclX@{}}
    \toprule
    Public family & $\Delta n$ & Status & Molecular role \\
    \midrule
    atom insert & $+1$ & core & Root or single-attachment atom birth \\
    atom delete & $-1$ & core & Legal atom death with connected successor \\
    atom restate & $0$ & core & Element--valence change at a retained site \\
    bond reorder & $0$ & core & Bond-order and local electronic change \\
    bond reroute & $0$ & core & Move an attached fragment to a legal target \\
    cycle close & $0$ & core & Add a legal bond that increases cycle rank \\
    cycle open & $0$ & core & Remove a cyclic bond while preserving connectivity \\
    ring-system restate & $0$ & accelerator & Coordinated cyclic electronic restatement \\
    \bottomrule
  \end{tabularx}
\end{table}

\paragraph{Birth and death.}
The factorized birth head supports a new root with zero existing neighbors and
a connected insertion with exactly one existing neighbor. It does not score an
atom attached simultaneously to two or more existing atoms. Consequently,
atom birth and death make the semantic process trans-dimensional, but the
operator set is not one-step inverse-closed for every deletion: the inverse of
deleting a degree-two bridging atom can require an unsupported multi-attachment
insertion. Such insertions are added only if a frozen reachability or path-cost
audit establishes a material task gap.

\paragraph{Connectivity and topology.}
Bond rerouting preserves an attached fragment while changing its attachment,
which can express positional and linker edits more directly than deleting and
rebuilding the fragment. Primitive cycle close and cycle open define the ring
topology support compositionally; neither depends on a finite catalog of named
ring templates. Ring-system restatement is retained as a coordinated
electronic accelerator for the bounded pilot. An exhaustive training-corpus
audit found that its stored teachers lower to valid sequences of cyclic bond
reorders, but the current primitive bond-reorder mask excludes every such
lowering. Removing it would therefore silently remove those teachers rather
than simplify an equivalent support.

Whole-ring growth and ring-system deletion are disabled in the bounded rebuild.
The former would reintroduce a finite topology catalog. The latter remains
excluded unless a frozen reachability audit establishes a meaningful
successor class or path benefit beyond primitive cycle opening. A future macro
can be added as an explicitly costed accelerator only if it preserves or
explicitly extends the declared reachable set and demonstrates a material
finite-budget gain.

The cycle namespaces differ across interfaces. Public \emph{cycle close} is
model family \texttt{cycle\_insert} and executor rule
\texttt{bond\_insert}; public \emph{cycle open} is model family
\texttt{cycle\_attach} and executor rule \texttt{bond\_delete}. The inherited
name \texttt{cycle\_attach} does not mean attaching a ring. Contracts and
result artifacts bind both names so that a direction cannot be silently
reversed.

\subsection{Rate network and action factorization}

The network encodes the current persistent-slot graph and progress variable,
then predicts a total marked-event hazard and a hierarchical conditional mark
law. At a high level,
\[
  p_\theta(a\mid x,t)
  =
  p_\theta(k\mid x,t)\,
  p_\theta(o,\eta\mid k,x,t),
\]
where $k$ is an operator family, $o$ denotes its concrete operands or matched
sites, and $\eta$ denotes typed payload such as an atom class or bond state.
Family-specific candidate builders and masks define the finite legal
coordinates scored by the second factor. Candidate scores are normalized only
over supported coordinates. Canonical-successor probabilities are obtained
after the complete factorized mark has been scored and executed.

The reference architecture uses a shared molecular message-passing encoder,
a pooled graph--time representation for the family gate, and family-specific
node, pair, or reroute scorers. The final width, depth, candidate features,
initialization, and parameter count are
\resultpending{EDITING_ARCHITECTURE_CONTRACT}. The property-agnostic base
network does not receive the immutable source molecule, optimization objective,
remaining budget, or protected mask. Those variables belong to the value or
controller, which permits all controllers to share one frozen base kernel.

\subsection{Endpoint programs and editing corpus}
\label{sec:editing-corpus}

\paragraph{Program compilation.}
Training records are executable paths between exact endpoint states. For a real
endpoint pair, the compiler tests direct semantic explanations, including
restate, reorder, reroute, birth/death, and cycle operations, before accepting a
delete-and-rebuild fallback. Every proposed path is replayed through the
production executor; its committed states, endpoint identity, charge policy,
operator sequence, aliases, and failure reason are recorded. Exact slot states,
not reparsed canonical strings, are the source of truth for replay.

\paragraph{Evidence lanes.}
The corpus contract distinguishes five data lanes:

\begin{enumerate}
  \item \texttt{observed\_local\_analogue}, for context-matched local
  substitutions and size changes between observed endpoints;
  \item \texttt{operator\_aware\_real\_endpoint}, which compiles observed
  endpoint pairs with direct Active8 operations before a delete--rebuild
  fallback;
  \item \texttt{linker\_positional\_topology\_analogue}, for two-cut linkers,
  positional isomers, attachment relocation, and topology-changing
  real-endpoint analogues;
  \item \texttt{real\_endpoint\_multistep\_path}, for short and medium
  compiler-generated paths between real endpoints with pair-relationship
  provenance retained; and
  \item \texttt{reversible\_synthetic\_walk}, for rare Active8 contexts,
  mixed-family composition, correction, and inverse coverage.
\end{enumerate}

Genuine observed-series provenance and observed action sequences are not
currently available. A path between real endpoints is therefore not described
as an observed series or an observed edit trajectory. Each record preserves
six evidence components separately: source endpoint, target endpoint, pair
relationship, path, intermediates, and action sequence. Real endpoints do not
make compiler-generated actions or executor-generated intermediates observed.
Synthetic paths establish executable support coverage; they are not evidence
of medicinal chemistry frequency. The frozen counts, unique sources and
scaffolds, path lengths, semantic-cell coverage, and effective training
coefficients are \resultpending{EDITING_CORPUS_CENSUS}.

\paragraph{Splits and sampling.}
The frozen split contract requires molecule, scaffold, authoritative
trace-membership groups, physically resolved pair-relationship groups, source
group, and transformation signature to be assigned before path construction;
no pair may cross partitions. The current census does not itself assign those
partitions. Because the previous final test aggregate was inspected during
diagnosis, the replacement model requires a new sealed final holdout. Training
samples hierarchically by data lane, series/scaffold/source group, semantic
capability cell, endpoint pair or path, and progress state. This prevents long
serializations, large relationship groups, or one common transformation from
dominating merely by row count. The final split artifact, mixture, and measured
effective coefficients are
\resultpending{EDITING_SAMPLER_CONTRACT}.

\subsection{Objective selection and no-waste training gates}
\label{sec:training-gates}

The editing experiments consume the productive embedded successor kernel, so
the replacement protocol compares the embedded successor likelihood in
\eqref{eq:productive-identity-loss} with the successor-generator objective in
\eqref{eq:successor-gm-loss}. Selected-mark training is retained as a bounded
ablation, not silently relabeled. The final objective, hazard weight, and
initialization are \resultpending{EDITING_OBJECTIVE_CONTRACT}.

The corpus contract remains
\texttt{DESIGN\_NOT\_TRAINING\_AUTHORIZED}. Authoritative trace-derived
membership-classifier receipts, physical relationship receipts, and the later
split-assignment artifact remain external prerequisites. A hash-valid schema
or a green unit-test suite cannot substitute for those physical inputs.

No long run is authorized directly from aggregate validation loss. A sequence
of immutable gates tests the intended scientific object:

\begin{enumerate}
  \item \textbf{D0/O0: corpus and support.} Every teacher executes, every
  successor fiber is recoverable, every capability cell is represented, and
  the operator basis meets frozen reachability and path-cost requirements.
  \item \textbf{T1: successor micro-overfit.} Unique-state panels must memorize
  their canonical successors; repeated-state panels must fit their empirical
  multi-successor law. Family mass alone is insufficient.
  \item \textbf{P50: gradient and collapse sentinel.} Fifty exact optimizer
  updates must expose every required family, preserve finite nonzero gradients,
  and avoid immediate capability collapse.
  \item \textbf{P500: mixed pilot.} Every required successor slice must improve
  over its registered baseline without deletion, size, or topology drift.
  \item \textbf{P2000: scientific pilot.} Held-analogue transport, cardinality
  adaptation, topology adaptation, calibration, and productive rollouts must
  pass before a full schedule is released.
\end{enumerate}

The optimizer, data-stream position, RNG state, and scheduler are continuous
across passing pilot stages; a pilot is not retrospectively warm-started with a
different law. The frozen full-run identity, selected checkpoint, and selection
metrics are \resultpending{EDITING_CHECKPOINT_FREEZE}.
